\mynote{Castoffs from the Introduction. Dunno if useful or not. Just keeping them here ATM}
\philnote{Doesn't make sense}
There is no denying that data creates opportunities, but it can also impose people its use. For example in extreme cases of \textit{datafication} in governance, where there is no need for engaging the public in any issue around the state because there is data that can simply be tallied up to enforce decisions. This means that all the intricacies, unique cultural peculiarities and imperfect organic ways people do things are discarded, and communities are boiled down to fabricated algorithmic identities.

\philnote{A bit of a clumily written paragraph.}
\petenote{Claiming stuff with no evidence. Heading to the loss of confidence in the author. style here, for me, it was the first thing I read, while ago now was that you sort of making quite a few claims. You weren't sort of referencing or evidencing or anything. And it began to feel a little bit like
the mind my confidence in the academic quality of it, or the, you know, the quality of the arguments}
But what happens to people in these communities? Do people just have to stay as subjects of this data or could this data be leveraged by them also to participate in these decision-making processes? There is a growing demand for data and evidence from the civil society when engaging in decision-making processes -- either at local or national level [REF]. Often there is lack of data for these groups to state their claims, but in some cases data is already generated. For example by the government conducting censuses and released the data as Governmental Open Data. However in many cases the value offers from that data are hidden from these community groups because of lack of access, know-how, skills or unevenly distributed resources. They are left to their own devices to make data work for them.

% Although the the amounts of data generated are vast, there is often paucity of information about how to use it in a purposeful way. 
% What is the role of data science and data scientist is it mealy a janitorial work for cleaning data

\philnote{Does corporatization of technology need further explanation?}
Data might be out there, however to condense it into the form that it is accessible and usable by everyone is not a straight forward task due do because of; lack of clear documentation and technical support; corporatization of technology; and absence of unified data standard across the board. Thesis expands on these issues further through reflection and discussion on the lessons learned from designing and building systems for accessible data in Chapter 4. However making data accessible is only the tip of the iceberg of effectively using it in civic advocacy and action. Chapter 5 reflects on both case studies to understand the linked processes around effective use of data and relevant social capital needed turn it into activities for local benefit. 

\mynote{Write here what has been successful, e.g., people still use SenseMyStreet}
Finally drawing on the literature the thesis reflects on the research conducted and the systems built for accessible data and their utility for communities hoping to instrument change, to propose a framework for effective use of data. Reflecting on case studies the framework helps to identify the key principles for successful engagement, contributing new understanding of the relationships between data, communities and knowledge networking and providing further design considerations for creators of technologies that aim to democratise data science practices use data for social good.

Chapter starts off with describing the role of data in the context of knowledge and value creation for civic action linked to complex decision-making processes. Moving away from the assumption that data is the `silver bullet', meaning if there is data information, knowledge and wisdom reveals itself without additional steps, naturally, automatically; to identifying the key principles of knowledge creation from data that leads to action. Chapter concludes with a suggested framework for effective use of data that helps understand the gaps and needed resources to promote data use within the civil society for actively participating in renewed democratic processes. 


so it is often natural for systems engineers and informaticians (i.e., data scientists) to bring the same methodology to accompany with the data.
This is also referred to as tacit knowledge which is hard or nearly impossible to communicate.
\cite{hakken_social_2003} also points to the shortcomings of expert systems not adequately dealing with the social, contextual and cultural dynamics at the centre of knowledge networking. By not recognising the social. We can’t ignore the social and just engineer or model intelligence, information and knowledge.

\petenote{weak defence for digital civics}
To respond to these challenges, new \textit{digital civics} \citep{olivier_digital_2015} approaches have emerged that recognise the importance of `social' in these processes. The affordance of internet of thing devices, and connected collaborative nature of the internet has enabled new models of knowledge production that promote participation and co-creation of information sources and services. Examples of such approaches are citizen sensing \citep{balestrini_city_2017,gabrys_programming_2014}, crowd sourcing for planning \citep{maskell_spokespeople:_2018,le_dantec_planning_2015} or even communities commissioning their own technologies to create a new data sources \citep{garbett_app_2016}. Projects like these have enabled citizens to generate alternative data sources to be used within the community to evidence issues and potentially challenge decisions. Although the data is now generated by the citizens themselves there is a considerable step from having the data and turning that data into actionable knowledge and activities for local benefit. This often involves working with multiple sets of data from different sources; accessing, cleaning, merging and interpreting them, and using visualisation techniques to tell a story - all which are activities that a technically literate and experienced data scientist would do.


\mynote{Castoffs from Chapter1}
This chapter describes related research around the use of data in knowledge creation and citizen participation in democracy. The use of data in civic life and within communities spans across a wide range of disciplines from technical fields like sensing networks, pervasive computing and information engineering to human computer interaction (HCI) and computer-supported collaborative work (CSCW). The first two sections of the literature review give an overview of the increasing use of data in decision-making through the lens of smart city and looks at the ways that citizens are involved or not involved in these data driver processes. Subsequent section describe the new ways of creating knowledge sources and opening them up for everyone to use and re-purpose. This also includes opportunities for new interpretations of open data sources published by the governments or other data generators. Fourth section describes the research around data production and different ways citizens can generate their own data sources, whether it is to contribute to science or to evidence local issues and contest decisions. Fifth section describes the current state-of the-art of data visualisation and storytelling through data which is essential for data science practices and the effective us of data. Drawing on work done by State of Open Data\footnote{\url{https://www.od4d.net/stateofopendata}} project and its contributors, section six tries to map out the relevant stakeholder groups around data production, use and dissemination of derived knowledge from it. It is essential to understand these complex relationships and roles of stakeholders, as we are increasingly coming to contend with significant challenges and questions relating to what data may be relevant and for whom; how to translate and represent it; as well as how different representations might be combined to better support communication and sense-making of complex information that sits at the boundaries between what may be private, shared and public. Building from that, section seven reviews the relevant literature about communities of practice and the knowledge networking within these communities. Finally the last section looks at the research into civic advocacy and action with the support of data and technology, including the existing models of successful and unsuccessful citizen decision-making. This literature review gives an overview of the underlying concepts touched in this thesis, however for each case study a more detailed and specific body of research is reviewed and presented.


\mynote{The following sections are things I took out from the literature review. They are either not needed, repetition of stuff, or things that are giving too much away and should probably go to other chapters/discussion. I don't really know yet!}

\mynote{Don't know yet where this fits in}
The promises from tech giants to `wire up' the city to solve some of the `wicked' problems might have been overstated or perhaps a bit premature. Many fancy dashboards and data silos later, society is starting to realise that data is not an end in itself. Having all the numbers does not give all the answers, to make sense and use the data, we must ask the right questions or possess the knowledge relevant to the issue at hand which can only obtained through a collaborative effort. This means we are not relaying on one piece of knowledge, but multiple and obtained from different sources -- smart city sensors, experts, city managers, researchers, decision-makers and citizens. There is an important role for designers and technologist to play here, working together with citizens to develop smart city tools for data collection, access, interpretation that promote civic participation and aim to democratise data science practices in the smart city. Similar design principles have guided the development of case studies in this thesis, where the developed tools and technology are direct responses to the needs of citizens.

\mynote{Don't know yet where this fits in}
All of these can be taken as context needed to produce and understand knowledge, because in these cases knowledge is always situated in activities engaged in the context of collective consciousness \citep{winograd_understanding_1986,hakken_social_2003}. In these situations knowledge mostly takes a tacit form \citep{polanyi_tacit_2009} which makes it difficult for individuals and organisations from outside the CoP to obtain any information. In order to understand the information it conveys, there is a need to understand what the context is and what does it mean \citep{hakken_social_2003}. At the same time this makes it difficult for people outside the domain to contribute to the knowledge base or identify additional resources needed that would benefit the practices.

In cyberspace people are increasingly relying on different data sources to derive knowledge. 
\philnote{?Who is pushed out here - needs clarification}
Often generated by somebody other than individuals involved in the CoPs and pushed out as they are, missing crucial context to really understand and use those data in knowledge creation. An example of this are datasets published as Open Government Data which represent the information collected for the purpose of governance and depict the viewpoint of the government. However that data is connected to organisational practices, which share some similarities with CoPs in terms of having a knowledge base in which it operates, but in reality it is inherently different from the organically created CoPs which do not possess the fixed hierarchical structures and tasked roles within (REF). This does not imply that that data is useless in community practices, it means that data needs to be contextualised by the community related to their practises. It is known from situation theory that information is in fact situated data and knowledge is information whose situation or context is known \citep{hakken_social_2003}. Case studies in this thesis explore ways that available data sources can be made into information that is useful for the communities and their practices. However, there should not be a confusion between \textit{information society} and \textit{knowledge society}. Whereas information society is more focused on generating and publishing data, knowledge society is more concerned with creating mechanisms and processes for effective use of those datasets and putting it into action towards improving the society as a whole. By designing tools with and for CoPs to collect data (SenseMyStreet) or help connect data to concerns (Data:In Place) enable people build stronger community networks, helping create and share knowledge for the community to act upon.

% contributing to an already existing systems (UO)
% difficult to achieve these links, also if they do not fit into the governmental structures they don't work. Data is just collected, but not acted upon
 On the one end it is difficult use data without knowing the context behind the generation and purpose. At the other end of the spectrum new tools are often creating new processes that do not link or interact with the existing ones, if they do not fit into the governmental structures they fail to produce any results. Data is just collected, but not acted upon. There is a lack of understanding of the hole process from data access to use and a clear framework to guide citizens in that process to use data in knowledge creation towards actionable results.

% dialog discourse
Although there are different focuses and strategies for civic advocacy, \citet{asad_tap_2017} emphasised that in complex civic processes, which usually involve more than one stakeholder, there is need to view issue advocacy across \textit{multiple sites} and venues. 


\philnote{What are the spaces in this context - don't follow.}
Developing a common language is critical in order to open spaces where common understandings of what data may be important for civic life can be developed for common purposes. Furthermore, a deeper examination of issues and tensions arising from processes of data classification, representation and interpretation becomes necessary in order to develop technologies and processes to better support collaborations around data capture, analysis, and sensemaking\citep{furnas_making_2005} , envisaged for a range of purposes from personal data management to community, civic and political engagement.

\philnote{We need more lining text from previous to this paragraph.}
% Look at the IBIS for HCI - Making Claims
% https://en.wikipedia.org/wiki/Issue-based_information_system
Trying to explore, map out and make sense of complex issues in a multi-stakeholder setting is not a new problem. It is often called dialogue mapping and associated with Issue Based Information Systems (IBIS). The original IBIS was developed in the early 1970s by Kunz and colleagues as a tool to support planning and policy design processes \citep{kunz_issues_1970} and with its computerised versions called graphical IBIS (gIBIS) following a decade later \citep{conklin_gibis:_1989,conklin_gibis:_1988}. 

% sentence about how now in addition to planners and government, people are added to the the mix of stakeholders
However, new challenges arise when trying to connect and navigate analytics and qualitative data, dialogues and opinions. There are two parts to this: (1) the visual sensemaking of analytics and (2) the exchange of opinions and judgements on the subject. Both are highly social and shared activities, where people collectively make sense of \citep{furnas_making_2005} analytics and then add their own layers of meaning to create knowledge.


\mynote{I don't know where this should go and even is it important.}
Finally, there is the question of what makes people participate and whether these new digital tools guarantee increase in participation? Multiple researchers have studied the way participation is affected when digital tools are implemented \citep{conroy_e-participation_2006,macintosh_e-democracy_2008,kim_e-participation_2012,macintosh_characterizing_2004,kitsing_online_2011}. Perhaps the most influential digital participation to governance is internet voting (I-voting) which is gaining popularity countries like Estonia, Switzerland, US, Canada and France. One of the pioneers in the area Estonia, has been using I-voting for all of its elections (local, parliamentary and EU parliamentary) since 2005. Although the participation through e-voting has increased, the overall voter turnout has stayed pretty much the same\footnote{\url{https://www.valimised.ee/en/archive/statistics-about-internet-voting-estonia}}. \citet{kitsing_online_2011} argues that though internet voting has a quicker turnover and is cheaper for the government, it does not necessarily make the process more inclusive. People who are already activity engaged will have another channel for participation, however the people who have never engaged will not necessarily start now because of a new technologically superior system. \citet{kitsing_online_2011} adds however, that having insights from a real world use case helps us understand these processes, and the heterogeneous nature of them, a bit better which opens up a space for reconfigurement of these processes.


\mynote{I don't know if this is something perhaps for Data:In Place chapter's, because there was citizens generating data that is not necessarily big data or numbers/stats}
Data however does not always have to refer to something that simply can be counted and represented as quantifiable values (e.g., `big data') \citep{michael_toward_2016,mcmillan_data_2016}, it can be much broader and diverse set of `things' often characterised as contemporary media such as 3D printed physical fabrications of data \citep{nissen_data-things:_2015}; travelling suitcases filled with stories and memories \citep{crivellaro_re-making_2016}; interactive physical graphs \citep{sweeney_pixel_2019}. Often a good way to engage people with data is through installations and art practices that use technology coupled with engineering practices to show data under a critical lens. Examples of these kind of projects are Take a Bullet for The City\footnote{\url{http://www.digiart21.org/art/take-a-bullet-for-the-city}}, which uses data from New Orleans Police Department to fire a blank shot every time when a report of a firearm discharge occurs and Invisible Airs\footnote{\url{http://yoha.co.uk/invisible}}, which uses different `machines' to visualise the public spending of Bristol City Council. In may cases this way of `releasing data' is far more effective than a flat machine readable file on some data portal and gets people more engaged with the issue in hand. In addition to the benefits of being able to combine digital and non-digital tools to visualise data, situated urban displays also enables people to keep the data in the place it was collected from and with the community it belongs to. For example using chalk and tape to visualise locally collected data to raise awareness around local issues and promote dialog between residents \citep{koeman_what_2014} or using public spaces to display private household data in order to promote behavioural change around energy usage \citep{moere_comparative_2011,bird_pulse_2010}. It is known however that the way in which different data comes to be interpreted is always entangled with and dependent upon both contexts and situations \citep{suchman_plans_1985,taylor_data--place:_2015}. While, there is a tendency for civic technology projects to place emphasis on quantitative data -- on the assumption that quantitative data may be relevant to people -- qualitative data describing things that cannot be easily quantified or measured are not less important. Indeed, it is often this type of `data' that helps people making sense of what matters and help generate actionable knowledge.
% How people make sense of data
% Chalk and tape, showing data in the community
% gameified interactions

\mynote{Castoffs of Chapter2}
These In the early stages of the investigation there 
Although issue of air quality is not the sole focus of the study around citizen participation in the smart city it is definitely one of the key drivers for the local authority and the citizens effected and involved in the study.

This phase provides an overview of activities for different groups and individuals who took part in the study. Purpose of the this step is to get an overview of the usage of the system and develop a theoretical basis for the analysis of the case study.

setting examples and committing individual efforts
Cycling community (SPACE for Gosforth) 14 July 2016 first met Ali
Ethnographic work and going to their events
Building relationships with the community
workshop - 22 participants
snowball sampling

initial interest group focus
ethnographic work - participating in meetings and in citizen organised events (setting up the group)
focus groups and interviews and participating in email exchange and in online media groups
Core group of 8 + me and Sean (core group was 5 people excluding me)
talk about the limitations of the sample

Demographics of the participants (retired ex professionals, professionals) health professionals also concerned about peoples health. Academia - planning, law, NHS professional, retired former council worker from a different council (city planner), general practitioner working for NHS, professional working in IT

Advocacy work outside - group rides, organising community meetings, raising awareness in the community (schools, churches, community events), active social media presents, participating in public consultations and schemes. 
insider outsider activities/online offline activities/ Part of networks

Regular cyclists for commuting, sport and leisure
Outside the group this analysis focuses on more people have signed up through the infrastructure ...
Also additional advocacy groups have been started to base on the model of this one
Having data to display and point at
 
The second part of the case study focuses on the evaluation of the use of the toolkit in the wild. The analysis is conducted through three phases: (1) analysis of the usage and activities on digital platforms built to support the community commissioning of sensors, (2) analysis of the discussions of a specific advocacy group on Facebook (i.e., community discussion), (3) analysis of audio recorded reflections of the core members from the same group on their advocacy efforts and the role of data generated using the toolkit (i.e., community reflection). Additionally, case study provides an analysis of observations of groups and individuals who are purposing the data in their advocacy efforts and the impact it has had. The data collected to report on the case study is however not definitive, as the toolkit is still operational, constantly evolving and being used by communities to gather data for exploring and evidencing issues in their neighbourhood.

and specific focus group that reflected on the advocacy effort and the role of generated sensor data through the toolkit. Research data collected in the process are in from of researchers notes, recorded focus group, documents produced by the advocacy group and discussions on social media platforms (i.e., Twitter and Facebook).


\petenote{This chapter seems to jump in at the deep end. I would reverse the two paragraphs. Start with something more like “This chapters explains…” and go onto to the overview of what sense my street is, the collaboration and so on.}

To investigate how data collected through a sensor commissioning toolkit can support citizen evidencing local concerns, and to understand the processes linked to operating a sustainable and transparent sensor commissioning toolkit, a number of deployments and engagement activities have been carried out in Newcastle Upon Tyne and surrounding areas over the period of three years. 

 SenseMyStreet toolkit was officially launched in the summer of 2017. Since then over 60 people have signed up, 7 groups of people have done environmental sensing with hand-held monitors loaned from the toolkit in their commutes to explore issues in their local area and three communities have commissioned stationary sensors to be deployed in their neighbourhood. Through constant engagement activities with the communities and feedback from online systems, the toolkit went through an iterative design process to involve into an operational citizen sensor commissioning toolkit.
 
 
 The data used in the analysis will include the fieldwork conducted by the author which studied different groups using the toolkit and appropriating the data gathered from the environmental monitors and the interactions recorded through the digital systems. 

The results from this project will be used to analyse the methods used in the design process and the overall challenges and opportunities associated with developing a transparent sensor commissioning toolkit, which allows people to collect, share, explore and purpose data collected using high precision environmental sensors.
% mention amount of sensor we have in the toolkit??


\mynote{Castoffs from Chapter3}
---Have you any information on new initiatives there has been made upon the data collection?---

When we started there’s was one established community/pressure group that worked on influencing the decisions about the built environment in their local area - http://spaceforgosforth.com/. As we started our work people have now organised in different areas in Newcastle and started their own local groups from the example of the first - https://www.spaceforheaton.com/,  http://www.spaceforjesmond.com/ and others. Aim for these groups is to engage their local area with the issues and activities that people can do change them — one of them drawing attention to the air quality and issues with traffic. Through these groups we have a wider range of engagement. Also they are all collaborating with each other and sharing their knowledge.

---Have you any information on CO2 levels or improvement of air quality?---

This is a more complicated question to answer. A side note that the main concern has shifted from CO2 to NO2 emissions from diesel cars which are very harmful to people’s health. UK government published their policy around tackling roadside NO2 in 26 July 2017 which requests a response from local governments (including Newcastle council) to respond with their local plan to reduce pollution in the city. UK government has been sued multiple times for not addressing the illegal levels of air pollution in different local authorities - https://www.clientearth.org/government-loses-third-air-pollution-case-judge-rules-air-pollution-plans-unlawful/. In Newcastle there’s 2 air quality management areas (https://www.newcastle.gov.uk/environment-and-waste/pollution/air-pollution/monitoring-air-quality) which are the main concern for the government. ATM the local council is working with the transportation providers, city planning officials, University (Including the Urban Observatory) and other experts to look at different solutions to improve the air quality in those areas mainly. However they are also trying to consult with the wider public around the issues of air quality, including the people who are actively involved with the SenseMyStreet project. The local government have to publish their proposed solution in March 2018. However a lot of people are worried because some of the other areas which are not in air quality management areas are still getting high values of pollution and if there’s going to be some change that will cause people to change their movement patterns the other areas can have a bigger impact. As we don’t know the plan at this moment I can’t speculate on what will happen. There needs to be a holistic approach to these things though including input from everyone - citizens, environmental sensors, experts, researchers and decision-makers (government). Also we need to look at things like social inclusion — sometimes making the environment better and air cleaner with clean air zones (taxing people in specific areas) will make people’s, who are from deprived areas, lives more difficult. 

However the only way to tackle the air quality issue is to go directly to the sources — ourselves. So it’s more about raising awareness in local communities and using the data to start conversations and draw people’s attention to the issue. The groups are working with council to deliver new schemes for active travel - https://www.newcastle.gov.uk/parking-roads-and-transport/cycling/streets-people-1, but the data is more useful for engaging the community and working on nudging people to change their behaviour.

---How are you communicating the project to the citizens and recruiting them?----

At first the we communicated information about the project in different community meetings which we attended and also a lot of local councils stakeholder meetings about the city. A lot of information is passed via word of mouth and people pointing others to the website or directly to us. Also I have worked with different community organisations and local interest groups for in the area for over three years now and have established a good relationships with them.

\section{Other initiatives related/Impact}

Planning proposals and changes (Streets for People, Cycle ambition fund)
Activist groups informing (Groups of concerned citizens) 
Community support
Sense Explorers (SenseMyStreets in Schools and Sean's stuff)

Press interest, SPAN, available resources
With the opening of Urban Sciences Building, my SenseMyStreet toolkit (case study 2) was featured in Newcastle Press10 , Chronicle Live11 and appeared as segment in BBC Look North (on 14/09/2017) lunchtime and evening show. Project was also discussed on blogs of communities taking part of the study. 

new innovations in smart cities
ncl.ac.uk/press/articles/archive/2017/09/sensemystreet 
chroniclelive.co.uk/news/north-east-news/smog-tyne-mine-what-newcastle-13610648  
spaceforheaton.com/articles
www.spaceforjesmond.com/2018/02/14/airquality\_20180411
ouseburntrust.org.uk/Handlers/Download.ashx?IDMF=1cd379c7-ec59-4d31-af11-6e1ce5e71508

The author advised the making and appeared in two BBC’s Geography: The Big Issues educational films made by Evans Woolfe Media, exploring the way that traffic affects the environment and people’s health. Films also featured the sensors in the SenseMyStreet toolkit

digital leaders smartest city of the year 2019